% GENERATED FILE. Edit canonical lessons, then run npm run generate:books.
% canonical-source-hash: fnv1a64:a765beb5a5df3528
% canonical-lessons: GU-C353-nagrik, GU-C353-janta, GU-C353-matdan, GU-C353-umedvar, GU-C353-cutavu

\chapter{Citizens and Votes}
\label{ch:gu-citizens-and-votes}
\hlchaptermodality{353}

\begin{chapteropening}
\textbf{By the end of this chapter:} \emph{I can talk about voting and candidates.}

Complete the last lesson of chapter 353: I can talk about voting and candidates.
\end{chapteropening}

\section[\texorpdfstring{\gu{નાગરિક} (nāgrik)}{નાગરિક (nāgrik)}]{\gu{નાગરિક} (nāgrik) — a citizen}

\label{lesson:GU-C353-nagrik}



\begin{quote}
Before the new one: say the Gujarati for the people of a state, then the Gujarati for to shape.
\end{quote}

\subsection*{You'll want to know: \gu{નાગરિક}}
\textbf{\gu{નાગરિક}} — \emph{nāgrik} — \textquotedblleft{}a citizen\textquotedblright{}.

\begin{quote}
\gu{દરેક} \gu{નાગરિકને} \gu{મત} \gu{આપવાનો} \gu{હક} \gu{છે}.
\end{quote}

\begin{quote}
\emph{darek nāgrikne mat āpvāno hak che.}
\end{quote}

\begin{quote}
Every citizen has the right to vote.
\end{quote}

A masculine or feminine noun, by the person meant.

\begin{cousinweb}
\textbf{\gu{નાગરિક}} (\emph{nāgrik}) is a tatsama, borrowed from Sanskrit \emph{nāgarika}, \textquotedblleft{}living in a town, civic\textquotedblright{} (Monier-Williams).
\end{cousinweb}

\begin{grammarlens}[title={What you've built}]
One more word for giving reasons and opinions.
\end{grammarlens}

\subsection*{Your turn}
\begin{itemize}
\raggedright
  \item \emph{Say it:} \emph{nāgrik}
  \item \emph{Say it:} \emph{nāgrik}, once more
  \item \emph{Say it:} \emph{darek nāgrikne mat āpvāno hak che.} — Every citizen has the right to vote.
  \item \emph{Recall:} say the Gujarati for the people of a state, then the Gujarati for to shape, then say \emph{nāgrik} again
\end{itemize}

\begin{tcolorbox}[breakable,colback=teal!4,colframe=teal!35!black,title={Before you move on}]
What does \gu{નાગરિક} mean? (\textquotedblleft{}A citizen\textquotedblright{}.) Say it once more.
\end{tcolorbox}
\section[\texorpdfstring{\gu{જનતા} (jantā)}{જનતા (jantā)}]{\gu{જનતા} (jantā) — the public}

\label{lesson:GU-C353-janta}



\begin{quote}
Before the new one: say the Gujarati for to shape, then the Gujarati for a citizen.
\end{quote}

\subsection*{You'll want to know: \gu{જનતા}}
\textbf{\gu{જનતા}} — \emph{jantā} — \textquotedblleft{}the public\textquotedblright{}.

\begin{quote}
\gu{જનતા} \gu{આ} \gu{નિયમથી} \gu{નારાજ} \gu{છે}.
\end{quote}

\begin{quote}
\emph{jantā ā niyamthī nārāj che.}
\end{quote}

\begin{quote}
The public is unhappy with this rule.
\end{quote}

A feminine noun.

\begin{cousinweb}
\textbf{\gu{જનતા}} (\emph{jantā}) is a tatsama, borrowed from Sanskrit \emph{janatā}, \textquotedblleft{}a body of people\textquotedblright{} (Monier-Williams).
\end{cousinweb}

\begin{grammarlens}[title={What you've built}]
One more word for giving reasons and opinions.
\end{grammarlens}

\subsection*{Your turn}
\begin{itemize}
\raggedright
  \item \emph{Say it:} \emph{jantā}
  \item \emph{Say it:} \emph{jantā}, once more
  \item \emph{Say it:} \emph{jantā ā niyamthī nārāj che.} — The public is unhappy with this rule.
  \item \emph{Recall:} say the Gujarati for to shape, then the Gujarati for a citizen, then say \emph{jantā} again
\end{itemize}

\begin{tcolorbox}[breakable,colback=teal!4,colframe=teal!35!black,title={Before you move on}]
What does \gu{જનતા} mean? (\textquotedblleft{}The public\textquotedblright{}.) Say it once more.
\end{tcolorbox}
\section[\texorpdfstring{\gu{મતદાન} (matdān)}{મતદાન (matdān)}]{\gu{મતદાન} (matdān) — voting}

\label{lesson:GU-C353-matdan}



\begin{quote}
Before the new one: say the Gujarati for a citizen, then the Gujarati for the public.
\end{quote}

\subsection*{You'll want to know: \gu{મતદાન}}
\textbf{\gu{મતદાન}} — \emph{matdān} — \textquotedblleft{}voting\textquotedblright{}.

\begin{quote}
\gu{રવિવારે} \gu{ગામમાં} \gu{મતદાન} \gu{થશે}.
\end{quote}

\begin{quote}
\emph{ravivāre gāmmā̃ matdān thaśe.}
\end{quote}

\begin{quote}
The village votes on Sunday.
\end{quote}

A neuter noun.

\begin{cousinweb}
\textbf{\gu{મતદાન}} (\emph{matdān}) is built from \textbf{\gu{મત}} (\emph{mat}), \textquotedblleft{}an opinion, a vote\textquotedblright{}, and Sanskrit \emph{dāna}, \textquotedblleft{}giving\textquotedblright{} (Monier-Williams): the giving of votes.
\end{cousinweb}

\begin{grammarlens}[title={What you've built}]
One more word for giving reasons and opinions.
\end{grammarlens}

\subsection*{Your turn}
\begin{itemize}
\raggedright
  \item \emph{Say it:} \emph{matdān}
  \item \emph{Say it:} \emph{matdān}, once more
  \item \emph{Say it:} \emph{ravivāre gāmmā̃ matdān thaśe.} — The village votes on Sunday.
  \item \emph{Recall:} say the Gujarati for a citizen, then the Gujarati for the public, then say \emph{matdān} again
\end{itemize}

\begin{tcolorbox}[breakable,colback=teal!4,colframe=teal!35!black,title={Before you move on}]
What does \gu{મતદાન} mean? (\textquotedblleft{}Voting\textquotedblright{}.) Say it once more.
\end{tcolorbox}
\section[\texorpdfstring{\gu{ઉમેદવાર} (umedvār)}{ઉમેદવાર (umedvār)}]{\gu{ઉમેદવાર} (umedvār) — a candidate}

\label{lesson:GU-C353-umedvar}



\begin{quote}
Before the new one: say the Gujarati for the public, then the Gujarati for voting.
\end{quote}

\subsection*{You'll want to know: \gu{ઉમેદવાર}}
\textbf{\gu{ઉમેદવાર}} — \emph{umedvār} — \textquotedblleft{}a candidate\textquotedblright{}.

\begin{quote}
\gu{આ} \gu{વખતે} \gu{ત્રણ} \gu{ઉમેદવાર} \gu{છે}.
\end{quote}

\begin{quote}
\emph{ā vakhte traṇ umedvār che.}
\end{quote}

\begin{quote}
There are three candidates this time.
\end{quote}

A masculine or feminine noun, by the person meant.

\begin{cousinweb}
\textbf{\gu{ઉમેદવાર}} (\emph{umedvār}) is usually traced to Persian \emph{umīdvār}, \textquotedblleft{}hopeful\textquotedblright{}, from \emph{umīd}, \textquotedblleft{}hope\textquotedblright{}: the one who hopes to be chosen.
\end{cousinweb}

\begin{grammarlens}[title={What you've built}]
One more word for giving reasons and opinions.
\end{grammarlens}

\subsection*{Your turn}
\begin{itemize}
\raggedright
  \item \emph{Say it:} \emph{umedvār}
  \item \emph{Say it:} \emph{umedvār}, once more
  \item \emph{Say it:} \emph{ā vakhte traṇ umedvār che.} — There are three candidates this time.
  \item \emph{Recall:} say the Gujarati for the public, then the Gujarati for voting, then say \emph{umedvār} again
\end{itemize}

\begin{tcolorbox}[breakable,colback=teal!4,colframe=teal!35!black,title={Before you move on}]
What does \gu{ઉમેદવાર} mean? (\textquotedblleft{}A candidate\textquotedblright{}.) Say it once more.
\end{tcolorbox}
\section[\texorpdfstring{\gu{ચૂંટાવું} (cū̃ṭāvũ)}{ચૂંટાવું (cū̃ṭāvũ)}]{\gu{ચૂંટાવું} (cū̃ṭāvũ) — to be elected}

\label{lesson:GU-C353-cutavu}



\begin{quote}
Before the new one: say the Gujarati for voting, then the Gujarati for a candidate.
\end{quote}

\subsection*{You'll want to know: \gu{ચૂંટાવું}}
\textbf{\gu{ચૂંટાવું}} — \emph{cū̃ṭāvũ} — \textquotedblleft{}to be elected\textquotedblright{}.

\begin{quote}
\gu{આ} \gu{વખતે} \gu{એક} \gu{જુવાન} \gu{ઉમેદવાર} \gu{ચૂંટાશે}.
\end{quote}

\begin{quote}
\emph{ā vakhte ek juvān umedvār cū̃ṭāśe.}
\end{quote}

\begin{quote}
This time a young candidate will be elected.
\end{quote}

Past (he …): \gu{તે} \gu{ઉમેદવાર} \gu{ચૂંટાયો}, \emph{te umedvār cū̃ṭāyo}.

\begin{cousinweb}
\textbf{\gu{ચૂંટાવું}} (\emph{cū̃ṭāvũ}) is the \emph{-āvũ} partner of \textbf{\gu{ચૂંટવું}} (\emph{cū̃ṭvũ}), \textquotedblleft{}to pick\textquotedblright{}, which CDIAL traces to Sanskrit \emph{cuṭáti}, \textquotedblleft{}cuts off\textquotedblright{} (CDIAL 4857): to be picked by the voters.
\end{cousinweb}

\begin{grammarlens}[title={What you've built}]
One more word for giving reasons and opinions.
\end{grammarlens}

\subsection*{Your turn}
\begin{itemize}
\raggedright
  \item \emph{Say it:} \emph{cū̃ṭāvũ}
  \item \emph{Say it:} \emph{cū̃ṭāvũ}, once more
  \item \emph{Say it:} \emph{ā vakhte ek juvān umedvār cū̃ṭāśe.} — This time a young candidate will be elected.
  \item \emph{Recall:} say the Gujarati for voting, then the Gujarati for a candidate, then say \emph{cū̃ṭāvũ} again
\end{itemize}

\begin{tcolorbox}[breakable,colback=teal!4,colframe=teal!35!black,title={Before you move on}]
What does \gu{ચૂંટાવું} mean? (\textquotedblleft{}To be elected\textquotedblright{}.) Say it once more.
\end{tcolorbox}
